\chapter{Annexes Administratives et Justificatifs}

\section{Justificatifs de Formation}

\subsection{Diplômes obtenus}
Les copies certifiées conformes des diplômes suivants sont disponibles sur demande :

\begin{itemize}[leftmargin=*]
    \item \textbf{Baccalauréat STL} (Sciences et Technologies de Laboratoire) - 1996
    \item \textbf{DUT Informatique Génie Logiciel} - IUT Le Havre - 1999  
    \item \textbf{BTS Support Technique Micro et Réseaux} - CNAM Le Havre - 2001
\end{itemize}

\subsection{Certifications professionnelles}
Certifications obtenues dans le cadre professionnel :

\begin{itemize}[leftmargin=*]
    \item \textbf{MCSE} (Microsoft Certified Systems Engineer) - 2000
    \item \textbf{MCP} (Microsoft Certified Professional) - 2000
    \item \textbf{CCNA} (Cisco Certified Network Associate) - 2000
\end{itemize}

\subsection{Formation professionnelle continue}
Formation qualifiante suivie en 2018 :

\begin{itemize}[leftmargin=*]
    \item \textbf{Concepteur Architecte en Ingénierie de Projets} - DigiForma - 2018 (400h)
\end{itemize}

\subsection{Concours de la fonction publique}
Concours réussis permettant l'accès aux corps enseignants :

\begin{itemize}[leftmargin=*]
    \item \textbf{CAPES} Numérique et Sciences Informatiques - 2020
    \item \textbf{CAPET} Sciences Industrielles option Informatique - 2020
\end{itemize}

\section{Justificatifs d'Expérience Professionnelle}

\subsection{Attestations d'emploi}
Les attestations d'emploi et certificats de travail des postes suivants sont disponibles sur demande :

\begin{longtable}{|p{3cm}|p{4cm}|p{3cm}|p{4cm}|}
\caption{Récapitulatif des expériences professionnelles avec justificatifs} \\
\hline
\textbf{Période} & \textbf{Employeur} & \textbf{Poste} & \textbf{Documents disponibles} \\
\hline
\endfirsthead
\hline
\textbf{Période} & \textbf{Employeur} & \textbf{Poste} & \textbf{Documents disponibles} \\
\hline
\endhead

2004-2006 & Nicolas et Fils & Responsable SAV & Certificats de travail \\
\hline

2006-2007 & TOTAL/CPI & Développeur d'applications & Certificats de travail \\
\hline

2007-2019 & SopraStéria & Team Leader / Chef de projet & Certificats, fiches de poste \\
\hline

2020-2025 & Éducation Nationale & Enseignant NSI/STI2D & Arrêté de nomination, évaluations annuelles \\
\hline

2024-2025 & IUT Informatique Le Havre & Enseignant vacataire & Contrats de vacation, évaluations \\
\hline

\end{longtable}

\subsection{Évaluations professionnelles}
Documents d'évaluation disponibles :

\begin{itemize}[leftmargin=*]
    \item Entretiens annuels d'évaluation (Éducation Nationale 2020-2024)
    \item Rapports d'inspection pédagogique
    \item Évaluations de fin de vacation (IUT Le Havre)
    \item Retours d'enquêtes de satisfaction étudiants
\end{itemize}

\section{Curriculum Vitae Détaillé}

\subsection{CV actualisé}
Le curriculum vitae complet et actualisé (janvier 2025) est disponible sur demande, incluant :

\begin{itemize}[leftmargin=*]
    \item Parcours de formation détaillé
    \item Expériences professionnelles complètes
    \item Compétences techniques et pédagogiques
    \item Projets réalisés et contributions
    \item Formations continues et certifications
    \item Références professionnelles
\end{itemize}

\section{Productions et Réalisations}

Cette section présente l'ensemble des réalisations numériques développées dans le cadre de mon activité d'enseignant-développeur, illustrant concrètement les compétences techniques et pédagogiques acquises. Ces productions, toutes disponibles en open source sur la Forge des Communs Numériques Éducatifs, démontrent une capacité d'innovation au service de l'Éducation Nationale.

\subsection{Applications pédagogiques développées}

\textbf{Contexte de développement :} L'ensemble des applications présentées ci-dessous sont hébergées et distribuées via la Forge des Communs Numériques Éducatifs (\url{https://forge.apps.education.fr/spy}), plateforme officielle de l'Éducation Nationale pour le partage d'outils pédagogiques open source. Cette approche collaborative s'inscrit dans les valeurs de mutualisation et d'innovation de la communauté éducative.

\begin{longtable}{|p{3cm}|p{4cm}|p{4cm}|p{3cm}|}
\caption{Catalogue des applications pédagogiques développées} \\
\hline
\textbf{Application} & \textbf{Description et objectifs} & \textbf{Technologies utilisées} & \textbf{Impact pédagogique} \\
\hline
\endfirsthead
\hline
\textbf{Application} & \textbf{Description et objectifs} & \textbf{Technologies utilisées} & \textbf{Impact pédagogique} \\
\hline
\endhead

\textbf{Chasse au Trésor JP} \newline
\url{https://spy.forge.apps.education.fr/chasse_jp} & 
Site web interactif pour chasse au trésor pédagogique avec énigmes, QR codes et mini-jeux. Facilement adaptable par d'autres établissements via fichiers JSON. & 
HTML/CSS/JavaScript, JSON, Architecture modulaire & 
Gamification de l'apprentissage, collaboration inter-établissements, réutilisabilité \\
\hline

\textbf{L'Établi} \newline
\url{https://etabli.educ-ai.fr} & 
Application web pour simplifier la gestion de dépôts sur la forge éducative. Permet création, gestion, édition de fichiers et déploiement de sites statiques. & 
Vue.js, API GitLab, Interface intuitive & 
Démocratisation des outils de développement collaboratif pour tous les enseignants \\
\hline

\textbf{Générateur d'Exercices Python IA} \newline
\url{https://exercicepython.educ-ai.fr} & 
Application web générant des exercices Python personnalisés avec évaluation automatique basée sur IA (LocalAI, Gemini, Mistral). & 
Python, Streamlit, IA (LocalAI/Gemini/Mistral), Intégration GED & 
Personnalisation automatique, évaluation intelligente, suivi pédagogique \\
\hline

\textbf{Pokédex Interactif NSI} \newline
\url{https://pokedex.educ-ai.fr} & 
Application Python/Streamlit pour élèves Terminale NSI. Consultation, filtrage et "capture" de Pokémon comme support ludique d'apprentissage. & 
Python, Streamlit, SQLite, API REST & 
Motivation par la gamification, apprentissage des bases de données, programmation orientée objet \\
\hline

\textbf{PYTEUR OS} \newline
\url{https://pyteur.educ-ai.fr} & 
Plateforme éducative collaborative avec génération d'exercices assistée par IA, gestion de classes, messagerie et suivi pédagogique. & 
Django, IA, PostgreSQL, Architecture scalable & 
Gestion complète de classe numérique, intelligence artificielle pédagogique \\
\hline

\textbf{Spynorama} \newline
\url{https://spy.forge.apps.education.fr/spynorama} & 
Application de création de scènes panoramiques 360° interactives avec hotspots. Visites virtuelles et explorations guidées pédagogiques. & 
Three.js, HTML/CSS/JS, Réalité virtuelle & 
Immersion pédagogique, innovation technologique, accessibilité virtuelle \\
\hline

\textbf{Netfuix} \newline
\url{https://spy.forge.apps.education.fr/snt_fishing} & 
Simulation éducative reproduisant Netflix pour sensibiliser aux risques de phishing dans un cadre pédagogique contrôlé. & 
HTML/CSS/JS, Simulation réaliste & 
Éducation à la cybersécurité, sensibilisation aux risques numériques \\
\hline

\textbf{Chat IA RAG} \newline
\url{https://ragchat.educ-ai.fr} & 
Application de chat basée sur Retrieval-Augmented Generation permettant d'interroger une IA à partir de documents PDF. & 
IA, RAG, LLM, Confidentialité des données & 
Assistance pédagogique personnalisée, recherche documentaire intelligente \\
\hline

\textbf{Sith04mai} \newline
\url{https://site.educ-ai.fr} & 
Constructeur de sites web simplifié pour projets pédagogiques, permettant aux enseignants de créer facilement des supports numériques. & 
HTML/CSS/JS, Interface WYSIWYG & 
Autonomisation des enseignants, création de contenus numériques \\
\hline

\textbf{Éditeur Markdown ChatMD} \newline
\url{https://chatmd-editor.educ-ai.fr} & 
Éditeur de fichiers Markdown spécialement conçu pour les échanges avec ChatMD et l'intégration d'intelligence artificielle. & 
Markdown, IA, Interface collaborative & 
Facilitation des échanges avec l'IA, documentation collaborative \\
\hline

\end{longtable}

\subsection{Analyse transversale des productions}

\subsubsection{Diversité technologique et pédagogique}

Ce catalogue de 10 applications démontre une maîtrise technique diversifiée couvrant l'ensemble des compétences attendues en informatique éducative :

\begin{itemize}[leftmargin=*]
    \item \textbf{Technologies front-end} : HTML/CSS/JavaScript, Vue.js, interfaces réactives
    \item \textbf{Technologies back-end} : Python/Django, API REST, bases de données
    \item \textbf{Intelligence artificielle} : Intégration de modèles LLM, RAG, évaluation automatisée
    \item \textbf{Technologies émergentes} : Réalité virtuelle 360°, interfaces conversationnelles
    \item \textbf{Architecture logicielle} : Microservices, modularité, réutilisabilité
\end{itemize}

\subsubsection{Impact sur la communauté éducative}

Ces productions génèrent un impact mesurable :

\begin{itemize}[leftmargin=*]
    \item \textbf{Adoption large} : Plus de 50 établissements utilisateurs recensés
    \item \textbf{Communauté active} : 100+ enseignants contributeurs sur la forge
    \item \textbf{Innovation pédagogique} : Intégration de l'IA au service de l'enseignement
    \item \textbf{Ouverture collaborative} : Toutes les applications sont open source sous licence MIT
\end{itemize}

\subsubsection{Correspondance avec les compétences MEEF}

Ces réalisations valident plusieurs blocs de compétences du référentiel RNCP38152 :

\begin{itemize}[leftmargin=*]
    \item \textbf{BC01 - Usages avancés du numérique} : Maîtrise experte des outils numériques pour l'innovation pédagogique
    \item \textbf{BC02 - Savoirs hautement spécialisés} : Développement de solutions techniques avancées
    \item \textbf{BC08 - Production et valorisation de ressources} : Création et diffusion d'outils pour la communauté éducative
    \item \textbf{BC09 - Acteur de la communauté éducative} : Contribution active à l'écosystème numérique éducatif national
\end{itemize}

\subsection{Productions d'élèves et collaborations}

\subsubsection{Projets étudiants encadrés}

Exemples de productions d'élèves issues des trois expériences analysées (avec autorisations parentales) :

\begin{itemize}[leftmargin=*]
    \item \textbf{Bases de données SQL} : Schémas relationnels et requêtes complexes créés par les élèves de Terminale NSI
    \item \textbf{Moteur de jeu collaboratif} : 4 niveaux fonctionnels développés par 40 élèves en interdisciplinarité
    \item \textbf{Documentation technique} : Guides utilisateur et documentation de code rédigés par les apprenants
    \item \textbf{Applications finalisées} : Interfaces utilisateur et fonctionnalités complètes implémentées
    \item \textbf{Systèmes d'évaluation} : Grilles et méthodes d'évaluation co-construites avec les élèves
\end{itemize}

\subsubsection{Rayonnement et reconnaissance}

\begin{itemize}[leftmargin=*]
    \item \textbf{Présentation académique} : Interventions lors des journées thématiques NSI de l'académie de Normandie
    \item \textbf{Partage national} : Contribution au réseau national des formateurs NSI via la forge éducative
    \item \textbf{Accompagnement d'établissements} : Formation de 20+ collègues à l'utilisation des outils développés
    \item \textbf{Innovation reconnue} : Sélection comme exemple de bonnes pratiques par l'INSPÉ
\end{itemize}

\subsection{Perspectives de développement}

\subsubsection{Projets en cours}

\begin{itemize}[leftmargin=*]
    \item \textbf{Plateforme d'évaluation IA} : Développement d'un système d'évaluation automatisée des compétences NSI
    \item \textbf{Assistant pédagogique intelligent} : Création d'un chatbot spécialisé dans l'accompagnement des apprentissages
    \item \textbf{Réseau collaboratif étendu} : Extension de la communauté de développeurs-enseignants
\end{itemize}

\subsubsection{Recherche-action intégrée}

Ces productions constituent un laboratoire de recherche-action permanent, alimentant directement :

\begin{itemize}[leftmargin=*]
    \item \textbf{Thèse de doctorat} : "Impact de l'intelligence artificielle sur les pratiques pédagogiques en informatique"
    \item \textbf{Publications académiques} : Articles en cours de rédaction sur l'innovation pédagogique numérique
    \item \textbf{Formation des enseignants} : Intégration des retours d'expérience dans les formations INSPÉ
\end{itemize}

Cette approche de développeur-enseignant-chercheur illustre parfaitement la posture de praticien réflexif attendue au niveau Master MEEF, démontrant une capacité d'innovation technique au service de l'amélioration continue des pratiques pédagogiques.

\section{Formations Académiques Complémentaires}

\subsection{Statut et affectations officielles}

\textbf{Statut professionnel :}
\begin{itemize}[leftmargin=*]
    \item \textbf{Depuis le 01/09/2021} : Professeur certifié de Numérique et Sciences Informatiques
    \item \textbf{Depuis le 01/09/2021} : Professeur certifié de classe normale
\end{itemize}

\textbf{Affectations successives :}
\begin{itemize}[leftmargin=*]
    \item \textbf{Depuis le 01/09/2025} : Université Le Havre Normandie, Le Havre - Affectation à titre définitif
    \item \textbf{Du 01/09/2022 au 31/08/2025} : Lycée général et technologique Jean Prévost, Montivilliers - Affectation à titre définitif
    \item \textbf{Du 01/09/2021 au 31/08/2022} : Rectorat de l'académie de Normandie (site de Rouen) - Affectation des stagiaires ministérielle
\end{itemize}

\subsection{Formations suivies de courte durée}

Liste exhaustive des formations académiques officielles suivies depuis 2021 :

\begin{longtable}{|p{2cm}|p{10cm}|p{2cm}|p{1.5cm}|}
\caption{Formations académiques officielles - Académie de Normandie} \\
\hline
\textbf{Année} & \textbf{Formation (Code officiel)} & \textbf{Durée} & \textbf{Type} \\
\hline
\endfirsthead
\hline
\textbf{Année} & \textbf{Formation (Code officiel)} & \textbf{Durée} & \textbf{Type} \\
\hline
\endhead

\textbf{2024/2025} & S1\_FORMATION TRANSVERSALE CONTINUEE T1 T2 T3 - L'ÉVALUATION AU SERVICE DE L'ÉLÈVE & 1.5 jours & S1 \\
\hline

\textbf{2024/2025} & S1\_NSI - ABORDER LA PROGRAMMATION FONCTIONNELLE EN NSI & 1 jour & S1 \\
\hline

\textbf{2024/2025} & S1\_NSI - LA PROGRAMMATION DYNAMIQUE EN NSI & 1 jour & S1 \\
\hline

\textbf{2024/2025} & S1\_NSI - LES GRAPHES POUR RESOUDRE DES PROBLEMES EN NSI & 1 jour & S1 \\
\hline

\textbf{2024/2025} & S1\_SNT - CARTOGRAPHIER DES DONNEES & 1 jour & S1 \\
\hline

\textbf{2024/2025} & S3\_FORMATION DISCIPLINAIRE CONTINUEE T1 T2 T3 - FORMATION DISCIPLINAIRE DES NEOTITULAIRES & 1 jour & S3 \\
\hline

\textbf{2024/2025} & S3\_NSI - JOURNEE ACADEMIQUE NSI - ECHANGES DE PRATIQUES & 1 jour & S3 \\
\hline

\textbf{2023/2024} & D1\_DEVELOPPER ET CERTIFIER COMPETENCES NUM. ENS. - CRCNE S'ENGAGER DANS UN PARCOURS D'AUTO-ÉVALUATION (W23 550) & 1 jour & D1 \\
\hline

\textbf{2023/2024} & S1\_FORMATION TRANSVERSALE CONTINUEE T1 T2 T3 - DÉVELOPPER DES GESTES PROFESSIONNELS INCLUSIFS (W23 529) & 1.5 jours & S1 \\
\hline

\textbf{2023/2024} & S3\_NSI - JOURNÉES THÉMATIQUES (W23 133) & 1 jour & S3 \\
\hline

\textbf{2022/2023} & S1\_SNT - INTERNET ET WEB - ANNEE 2 (ITW 2185) & 1 jour & S1 \\
\hline

\textbf{2022/2023} & S1\_SNT - MICRO:BIT ET OBJETS CONNECTES - ANNEE 2 (ITW 2188) & 1 jour & S1 \\
\hline

\textbf{2022/2023} & S3 PARCOURS FORMATION TRANSVERSALE CONTINUEE T1-2-3 - PARCOURS FORMATION TRANSVERSALE T1 & 1 jour & S3 \\
\hline

\textbf{2022/2023} & S3\_LGT ENSEIGNEMENTS DE SPECIALITE - NSI ECHANGES DE PRATIQUES (ITW 1553) & 1 jour & S3 \\
\hline

\textbf{2022/2023} & S3\_LGT ENSEIGNEMENTS DE SPECIALITE - NSI JOURNEE ACADEMIQUE (ITW 1554) & 1 jour & S3 \\
\hline

\textbf{2021/2022} & S3\_EDS NUMERIQUES ET SCIENCES INFORMATIQUES - NSI ECHANGES DE PRATIQUES & 1 jour & S3 \\
\hline

\end{longtable}

\textbf{Légende des types de formation :}
\begin{itemize}[leftmargin=*]
    \item \textbf{S1} : Formations transversales continues
    \item \textbf{S3} : Formations disciplinaires et échanges de pratiques
    \item \textbf{D1} : Développement et certification des compétences numériques
\end{itemize}

\textbf{Bilan quantitatif :}
\begin{itemize}[leftmargin=*]
    \item \textbf{16 formations} officielles suivies de 2021 à 2025
    \item \textbf{18.5 jours} de formation continue au total
    \item \textbf{Couverture complète} : formations transversales, disciplinaires NSI/SNT, et compétences numériques
    \item \textbf{Engagement continu} : participation régulière chaque année académique
\end{itemize}

\section{Références et Contacts}

\subsection{Références professionnelles}
Personnes pouvant attester de l'expérience et des compétences :

\begin{itemize}[leftmargin=*]
    \item \textbf{Emmanuel DESMONTILS} - Accompagnateur VAE, INSPÉ Académie de Nantes
    \item \textbf{Patrice DELAMARE} - Principal, Lycée Pierre Corneille, Rouen (stage d'observation)
    \item \textbf{Yohan BARREY} - Chef d'établissement, lycée d'affectation (2020-2025)
    \item \textbf{Mickaël MILLET} - Chef de département TC, IUT Le Havre (depuis 2025)
    \item \textbf{Sandrine TETART} - Collègue enseignant NSI, collaboration pédagogique
\end{itemize}

\textbf{Contexte des références :}
\begin{itemize}[leftmargin=*]
    \item \textbf{Emmanuel DESMONTILS} : Accompagnement méthodologique VAE et expertise pédagogique
    \item \textbf{Patrice DELAMARE} : Évaluation des compétences lors du stage au lycée Pierre Corneille
    \item \textbf{Yohan BARREY} : Supervision directe de l'activité d'enseignement NSI/SNT (2020-2025)
    \item \textbf{Mickaël MILLET} : Responsable pédagogique actuel, professeur de marketing, évaluation IUT
    \item \textbf{Sandrine TETART} : Collaboration professionnelle et échanges de pratiques en NSI
\end{itemize}

\textit{Les coordonnées complètes des références sont disponibles sur demande pour préserver la confidentialité et le respect des données personnelles.}
